\chapter{多任务微调数据集构建与清洗}\label{chap:dataset}

\section{引言}\label{sec:dataset-intro}

\cref{chap:reinforce}基于TQC算法与HER机制，为Reach、Push、Slide与PickAndPlace四项机械臂操作任务分别训练了单任务专家策略。但这些专家策略以MLP网络实现，输入与输出均为低维连续浮点数向量，大语言模型无法直接处理。要将这些策略用于\cref{chap:finetune}中的LLM机械臂控制，还需将物理观测轨迹转化为结构化文本，同时保证转化后的动作序列在物理引擎中可复现。

为此，本章围绕多任务微调数据集的构建与清洗展开。\cref{sec:data-transformation-design}设计从物理状态到自然语言的异构数据转化机制，涵盖观测空间的语义映射与连续数值的离散化表示，并定义三角色对话模板。\cref{sec:trajectory-collection}提出基于物理引擎闭环回放的轨迹采集与清洗方法，通过逐步误差追踪与环境自身任务成功判定来验证截断后动作序列的复现性。\cref{sec:data-quality-analysis}从匹配率、漂移误差与逐步误差三个方面对清洗结果进行定量分析。

\section{物理状态到自然语言的异构数据转化}\label{sec:data-transformation-design}
\subsection{数据转化设计原则与方法}\label{sec:data-transformation}

LLMs接收离散的自然语言文本序列，而\cref{sec:expert-training}中专家策略所运行的MuJoCo仿真环境，观测空间与动作空间均为连续浮点数向量。两者之间存在直接的格式差异：仿真环境输出的扁平数值数组缺乏结构，无法被语言模型的分词器有效编码。为此，本节设计了一套异构数据转化机制，将物理观测轨迹转化为带语义标注的结构化文本。该转化机制包含观测空间的语义映射与连续数值的离散化表示两部分。

在观测空间的语义映射方面，本研究将\cref{sec:rl-environment}中定义的各任务观测向量按物理含义切分为语义字段，为每个字段附加变量名与物理含义注释。不同任务的观测维度不同，因此各任务分别定义了映射规则，如\cref{tab:observation-mapping}所示。Reach任务的状态空间为10维，仅包含末端执行器与夹爪的运动学信息；Push、Slide与PickAndPlace任务由于涉及目标物块操作，观测空间扩展至25维，增加了物块的位姿、相对运动学信息及角速度。各任务观测字段的完整定义见\cref{chap:detail-env}。

\begin{table}[htbp]
\centering
\caption{各任务观测状态空间的语义切分方案}\label{tab:observation-mapping}
\begin{tabular}{llc}
\toprule
语义字段 & 物理含义 & 维度 \\
\midrule
\multicolumn{3}{l}{\textit{基础观测(全部任务共有)}} \\
gripper\_pos & 末端执行器三维位置 & 3 \\
gripper\_vel & 末端执行器三维线速度 & 3 \\
gripper\_open & 夹爪左右指间相对距离 & 2 \\
gripper\_open\_vel & 夹爪开合速度 & 2 \\
\midrule
\multicolumn{3}{l}{\textit{扩展观测(Push/Slide/PickAndPlace任务)}} \\
object\_pose & 目标物块全局位姿(位置与姿态) & 6 \\
object\_rel\_posvel & 物块相对末端执行器的位置与速度 & 6 \\
object\_angular\_vel & 物块全局角速度 & 3 \\
\midrule
观测维度合计 & & Reach: 10 / 其他: 25 \\
\bottomrule
\end{tabular}
\end{table}

在连续数值的离散化表示方面，本研究对所有物理状态参数和动作参数统一保留至小数点后4位，并以Python列表的字符串形式编码。精度选择需要在两方面之间权衡：保留位数过少会损失精度，降低动作向量的控制准确性；保留位数过多则会增加单条样本的Token序列长度，降低训练效率。实际采用的4位小数精度足以表征亚毫米级的位置差异，达到两者间的平衡。

\subsection{对话模板设计与数据封装}\label{sec:dialogue-template}

完成观测空间的语义映射后，还需将结构化文本封装为大语言模型可接受的对话格式。本研究采用三角色对话模板，由系统指令(System)、用户指令(User)和辅助指令(Assistant)组成，各角色的内容与作用如\cref{tab:prompt-template}所示。

\begin{table}[htbp]
\centering
\caption{三角色对话模板的职责与内容定义}\label{tab:prompt-template}
\begin{tabular}{lp{5.4cm}p{5.4cm}}
\toprule
角色 & 内容组成 & 设计意图 \\
\midrule
System & 任务角色设定；包含动作空间先验约束 & 为模型提供身份定义与输出规范 \\
User & 当前观测状态语义表述；目标完成信息 & 提供完整的决策上下文 \\
Assistant & 4维连续动作向量 & 去除冗余文本，专注数值映射学习 \\
\bottomrule
\end{tabular}
\end{table}

\cref{fig:prompt-example}展示了Reach任务的一个完整对话模板实例。System Prompt将模型设定为7自由度Fetch移动机械臂的专家控制器，指定控制目标为末端执行器的位置到达任务，并将输出动作空间约束为4维连续变量$(\text{dx}, \text{dy}, \text{dz}, \text{gripper\_control})$，各维度数值严格限制在$[-1.0, 1.0]$区间内。

User Prompt将环境输入分为当前观测状态与目标完成情况两部分。当前观测状态包含\cref{sec:data-transformation}所定义的末端执行器位置、速度及夹持器状态等物理状态变量，各数值均保留4位小数精度；目标完成信息包含当前已达到状态与期望目标状态的对比。

\begin{figure}[htbp]
  \centering
  \includegraphics[width=0.60\textwidth]{figures/prompt_example.png}
  \caption{Reach任务的三角色对话模板示例}
  \label{fig:prompt-example}
\end{figure}

Assistant Prompt是模型需要学习的目标输出。训练阶段，其内容来自\cref{sec:expert-training}训练的MLP专家网络，同样以4位小数精度表示。为减少干扰，动作指令仅保留4维连续动作的纯数值部分，不含文本描述与格式化符号。推理阶段，Assistant Prompt对模型不可见，模型自主生成动作向量后解析为numpy数组送入仿真环境执行。

在数据封装层面，上述多轮交互数据统一序列化为JSONL格式存储。由于后续实验将基于多个不同架构的开源模型进行微调，数据构建阶段不显式嵌入任何特定模型的特殊标记，保持格式与模型无关。待进入监督式微调(SFT)阶段时，再根据所选基座模型的分词器规则动态注入对应的格式化Token，使同一份数据集可在不同架构间复用。

\section{闭环轨迹采集与数据清洗机制}\label{sec:trajectory-collection}

在构建供大语言模型微调使用的专家数据集时，数据质量将直接影响下游控制模型的性能。\cref{sec:data-transformation}中的异构数据转化对连续数值做了截断处理，截断引入的浮点误差在动作序列逐帧执行中可能逐步累积，使原始轨迹在物理引擎中无法被精确复现。本节针对这一问题设计采集、回放与清洗的数据构建流程。

\begin{figure}[htbp]
\centering
\includegraphics[width=0.55\textwidth]{figures/dataset_workflow_collect}
\caption{闭环轨迹采集阶段流程}\label{fig:trajectory-collect}
\end{figure}

在如\cref{fig:trajectory-collect}所示的轨迹采集阶段，本研究调用\cref{sec:expert-training}训练的各任务MLP专家网络直接在环境中闭环运行。专家网络根据当前观测输出动作指令，环境执行该动作后返回新观测。轨迹数据以回合为基本单位记录，每条轨迹包含从初始状态到任务成功终止的完整状态-动作序列。采集时仅保留成功完成任务的回合，失败则丢弃并重采。每项任务各采集500条成功轨迹。采集结束后，对连续数值进行小数截断，再按\cref{sec:dialogue-template}定义的三角色对话模板进行语义化封装，生成JSONL格式数据。

\begin{figure}[htbp]
\centering
\includegraphics[width=0.50\textwidth]{figures/dataset_workflow_replay}
\caption{物理引擎回放验证阶段流程}\label{fig:trajectory-replay}
\end{figure}

为验证截断后数据的可复现性，本研究设计了如图\cref{fig:trajectory-replay}所示的闭环回放验证机制。对每条语义化轨迹，系统先将其中的动作向量从文本格式逆向解析为数值序列，再将MuJoCo仿真环境重置为该轨迹记录的初始状态，重置内容包括目标位置与物块初始坐标，然后按截断后的动作序列逐帧执行，过程中不做实时推理。回放过程中，系统同步计算每步的环境实际状态与记录状态之间的逐步偏差，即该步目标物体实际位置与记录值之间的欧氏距离。回放结束后计算最终漂移误差，即回放终止时实际状态与记录终止状态之间的位置偏差。

在\cref{fig:trajectory-clean}所示的清洗阶段，系统通过环境内置的\texttt{is\_success}判定函数评估任务是否达成。仅当回放后任务仍判定为成功时，该轨迹才被保留；否则说明截断后的动作序列不足以复现原始效果，予以剔除。清洗后各任务的数据质量分析见\cref{sec:data-quality-analysis}。

\begin{figure}[htbp]
\centering
\includegraphics[width=0.55\textwidth]{figures/dataset_workflow_clean}
\caption{数据清洗与输出阶段流程}\label{fig:trajectory-clean}
\end{figure}

\section{专家数据集质量分析}\label{sec:data-quality-analysis}

经过\cref{sec:trajectory-collection}所述的物理引擎回放验证与清洗后，各任务的轨迹匹配率与误差统计如\cref{tab:trajectory-removal-rates}所示。四项任务在采集阶段各产生500条成功轨迹，但清洗后的留存率差异明显，反映了各任务物理动力学特性的不同。

Reach任务仅涉及机械臂末端的运动学控制，不含物体接触，截断误差只影响末端执行器的位置精度。该任务的平均逐步误差为0.005m，最终漂移误差接近零，500条轨迹全部通过复现验证，匹配率为100\%。

Push与PickAndPlace任务涉及刚体接触与抓取操作，接触力使逐步误差在序列执行中有所累积，平均逐步误差分别达到0.016m与0.022m，最终漂移误差分别为0.018m与0.024m。但这两类任务的接触过程相对稳定，刚体约束限制了截断误差的进一步累积，匹配率仍分别保持在90.6\%与92.6\%。

Slide任务的误差放大最为明显。该任务要求机械臂通过瞬时碰撞将物块击打至低摩擦桌面上的目标位置，过程对初始接触状态与冲量十分敏感。浮点截断误差在碰撞与长距离滑行的动力学积分中被放大，平均逐步误差达到0.032m，最终漂移误差升至0.050m，匹配率仅为70.0\%，30\%的轨迹因复现失败而被剔除。

\begin{table}[htbp]
\centering
\caption{各任务轨迹回放清洗结果与误差统计}\label{tab:trajectory-removal-rates}
\begin{tabular}{lccc}
\toprule
任务 & 匹配率 & 平均漂移误差/m & 平均逐步误差/m \\
\midrule
Reach & 100.0\% & 0.001 & 0.005 \\
Push & 90.6\% & 0.018 & 0.016 \\
PickAndPlace & 92.6\% & 0.024 & 0.022 \\
Slide & 70.0\% & 0.050 & 0.032 \\
\bottomrule
\end{tabular}
\end{table}

通过物理复现筛查，本研究剔除了对初始条件敏感及含累积误差的不稳定轨迹，留存轨迹均满足状态-动作因果一致性要求，可作为后续大语言模型指令微调的数据来源。